% Luc Destombe --- licence CC BY-NC-SA 4.0
% https://creativecommons.org/licenses/by-nc-sa/4.0/deed.fr
% Fichier autonome : compiler avec pdflatex, deux fois (signets, nombre de pages).
\documentclass[11pt,a4paper]{article}
\makeatletter
\newif\iflycee@unecolonne
\lycee@unecolonnetrue
\newif\iflycee@palatino
\lycee@palatinotrue

\RequirePackage[utf8]{inputenc}
\RequirePackage[T1]{fontenc}
\RequirePackage[french]{babel}
\RequirePackage{etoolbox}

\newcommand{\ShorthandsOff}{\@ifundefined{shorthandoff}{}{\shorthandoff{;:!?}}}
\newcommand{\ShorthandsOn}{\@ifundefined{shorthandon}{}{\shorthandon{;:!?}}}
\ShorthandsOff
\AtBeginDocument{\ShorthandsOn}
\AtBeginEnvironment{tikzpicture}{\ShorthandsOff}

\RequirePackage{geometry}
\RequirePackage{fancyhdr}
\RequirePackage{lastpage}
\RequirePackage{multicol}
\RequirePackage{array}
\RequirePackage{enumitem}
\RequirePackage{xcolor}
\RequirePackage{needspace}

\RequirePackage{amsmath,amssymb}
\RequirePackage{mathpazo}
\DeclareMathSizes{10.95}{10.95}{8}{5.5}
\begingroup\catcode`\_=\active \catcode`\^=\active
\gdef_#1{\sb{\scriptscriptstyle #1}}
\gdef^#1{\sp{\scriptscriptstyle\lycee@moinsepais #1}}
\endgroup
\newcommand\lycee@moins{\mathbin{%
  \setbox\z@\hbox{$\m@th\scriptscriptstyle\lycee@moinsorig$}%
  \dimen@=.55\wd\z@ \advance\dimen@-.5pt
  \hbox to.75\wd\z@{\hss$\m@th\scriptscriptstyle\vcenter{\hbox{%
    \tikz\draw[line width=.5pt,line cap=round](0,0)--(\the\dimen@,0);}}$\hss}}}
\begingroup\lccode`\~=`\- \lowercase{\endgroup\let~}\lycee@moins
\newcommand\lycee@moinsepais{\mathcode`\-="8000 }
\AtBeginDocument{\mathchardef\lycee@moinsorig=\mathcode`\-\relax}
\AtBeginDocument{\catcode`\_=12 \mathcode`\_="8000
  \catcode`\^=12 \mathcode`\^="8000 }
\newcommand\lycee@exposants{%
  \fontdimen13\textfont2=.48\fontdimen6\textfont2
  \fontdimen14\textfont2=.48\fontdimen6\textfont2
  \fontdimen15\textfont2=.48\fontdimen6\textfont2
  \fontdimen13\scriptfont2=.48\fontdimen6\scriptfont2
  \fontdimen14\scriptfont2=.48\fontdimen6\scriptfont2
  \fontdimen15\scriptfont2=.48\fontdimen6\scriptfont2}
\every@math@size\expandafter{\the\every@math@size\lycee@exposants}
\newlength{\retraitcalculs}
\setlength{\retraitcalculs}{2em}

\RequirePackage{tikz}
\usetikzlibrary{arrows.meta,calc,babel,shapes.misc}
\RequirePackage[most]{tcolorbox}
\tcbuselibrary{skins,breakable}

\definecolor{coulDef}{HTML}{1F4E79}
\definecolor{coulProp}{HTML}{7030A0}
\definecolor{coulRem}{HTML}{C55A11}
\definecolor{coulEx}{HTML}{2E7D32}
\definecolor{coulExo}{HTML}{B03A2E}
\definecolor{coulPart}{HTML}{1F4E79}
\definecolor{coulMeth}{HTML}{1B6E4F}
\definecolor{coulCourbe}{HTML}{0B5ED7}
\definecolor{coulCourbeB}{HTML}{C00000}
\definecolor{coulCourbeC}{HTML}{2E7D32}
\definecolor{coulGrille}{HTML}{8C8C8C}
\definecolor{coulRep}{HTML}{1B6E4F}
\definecolor{coulBar}{HTML}{2E7D32}

\newcommand{\R}{\mathbb{R}}
\newcommand{\Rs}{\mathbb{R}^{*}}
\renewcommand{\le}{\leqslant}
\renewcommand{\ge}{\geqslant}
\newcommand{\vect}[1]{\overrightarrow{#1}}
\newcommand{\norme}[1]{\left\lVert #1 \right\rVert}
\newcommand{\intervalle}[2]{\left[#1\,;#2\right]}
\RequirePackage{cancel}
\renewcommand{\CancelColor}{\color{coulExo}}
\newcommand{\fois}[1]{\times{\color{coulExo}#1}}
\newcommand{\barre}[1]{\cancel{#1}}

\tikzset{grille/.style={coulGrille,line width=0.4pt}}
\newcommand{\quadrillage}[4]{\draw[grille,step=1] (#1,#2) grid (#3,#4);}
\tikzset{
  croix/.style={cross out,draw,line width=0.6pt,inner sep=0pt,minimum size=4pt},
  croixdroite/.style={croix,rotate=45,line width=0.8pt},
}
\newcommand{\pt}[3]{\path (#1) node[croix]{} node[#2,font=\small,inner sep=2.5pt]{$#3$};}

\newcommand{\lycee@repere}[4]{%
  \draw[grille,step=1] (#1,#2) grid (#3,#4);
  \draw[-{Stealth[length=2mm]},thick] (#1,0)--({#3+0.4},0) node[below left=-1pt]{$x$};
  \draw[-{Stealth[length=2mm]},thick] (0,#2)--(0,{#4+0.4}) node[below left=-1pt]{$y$};
  \node[below left,font=\scriptsize,inner sep=1.5pt] at (0,0) {$0$};}
\newcommand{\lycee@gradx}[1]{\foreach \k/\l in {#1}{%
  \draw (\k,0.07)--(\k,-0.07) node[below,font=\scriptsize,inner sep=1.5pt]{$\l$};}}
\newcommand{\lycee@grady}[1]{\foreach \k/\l in {#1}{%
  \draw (0.07,\k)--(-0.07,\k) node[left,font=\scriptsize,inner sep=1.5pt]{$\l$};}}
\newcommand{\lycee@intff}[2]{\left[#1\,;#2\right]}
\newcommand{\lycee@intfo}[2]{\left[#1\,;#2\right[}
\newcommand{\lycee@intof}[2]{\left]#1\,;#2\right]}
\newcommand{\lycee@intoo}[2]{\left]#1\,;#2\right[}
\newcommand{\lycee@ens}[1]{\left\{#1\right\}}
\AtBeginDocument{%
  \providecommand{\repere}{\lycee@repere}%
  \providecommand{\gradx}{\lycee@gradx}%
  \providecommand{\grady}{\lycee@grady}%
  \providecommand{\Cf}{\mathcal{C}\sb{\scriptscriptstyle f}}%
  \providecommand{\Cg}{\mathcal{C}\sb{\scriptscriptstyle g}}%
  \providecommand{\intff}{\lycee@intff}%
  \providecommand{\intfo}{\lycee@intfo}%
  \providecommand{\intof}{\lycee@intof}%
  \providecommand{\intoo}{\lycee@intoo}%
  \providecommand{\ens}{\lycee@ens}}

\newlist{questions}{enumerate}{3}
\setlist[questions,1]{label=\arabic*\textdegree),leftmargin=*,itemsep=2pt,topsep=3pt}
\setlist[questions,2]{label=\alph*),leftmargin=*,itemsep=1pt,topsep=2pt}
\setlist[questions,3]{label=\roman*),leftmargin=*,itemsep=1pt,topsep=1pt}
\setlist[itemize]{label=\textbullet,leftmargin=*,itemsep=2pt,topsep=3pt}

\newcommand{\besoinplace}[1]{\par\penalty\@M\begingroup
  \dimen@=#1\relax\dimen@ii=\pagegoal\advance\dimen@ii-\pagetotal
  \ifdim\dimen@>\dimen@ii\ifdim\dimen@ii>\z@\vfil\fi\break\fi\endgroup}

\newcommand{\lycee@pied}{\thepage/\pageref{LastPage}}
\newcommand{\entete}[2]{%
  \pagestyle{fancy}\fancyhf{}%
  \fancyhead[L]{\footnotesize\color{coulPart}\textbf{#1}}%
  \fancyhead[R]{\footnotesize\color{coulPart}\textbf{#2}}%
  \fancyfoot[C]{\footnotesize\lycee@pied}%
  \renewcommand{\headrulewidth}{0.4pt}%
  \renewcommand{\headrule}{\hbox to\headwidth{%
    \color{coulPart!40}\leaders\hrule height \headrulewidth\hfill}}}

\RequirePackage{ccicons}
\newcommand{\lycee@licence}{{\scriptsize\color{gray}%
  \copyright\ Luc Destombe --- \ccbyncsa\ CC BY-NC-SA 4.0}}
\newif\iflycee@licence \lycee@licencetrue
\newcommand{\sanslicence}{\lycee@licencefalse}
\AtBeginDocument{\iflycee@licence\AddToHookNext{shipout/foreground}{%
  \put(0,0){\rlap{\kern\dimexpr1in+\hoffset+\oddsidemargin+\textwidth\relax
    \llap{\raisebox{-\dimexpr1in+\voffset+\topmargin+\headheight+\headsep
      +\textheight+\footskip\relax}{\lycee@licence}}}}}\fi}

\newcommand{\formulesserrees}{%
  \setlength{\abovedisplayskip}{3pt}\setlength{\belowdisplayskip}{3pt}%
  \setlength{\abovedisplayshortskip}{2pt}\setlength{\belowdisplayshortskip}{3pt}}

\setlength{\parindent}{0pt}

\RequirePackage[implicit=false,hidelinks,pdfencoding=auto,psdextra,bookmarksopen,
  bookmarksopenlevel=1,pdfcreator={},pdfproducer={}]{hyperref}
\RequirePackage{bookmark}
\hypersetup{pdfauthor={Luc Destombe},pdfkeywords={CC BY-NC-SA 4.0}}
\newcounter{lycee@signet}
\newcommand{\signet}[3][]{\stepcounter{lycee@signet}%
  \edef\lycee@dest{\if\relax\detokenize{#1}\relax
    signet.\arabic{lycee@signet}\else#1\fi}%
  \hypertarget{\lycee@dest}{}\bookmark[level=#2,dest=\lycee@dest]{#3}}
\pdfstringdefDisableCommands{%
  \def\R{R}\def\vect#1{#1}\def\Cf{Cf}\def\Cg{Cg}\def\fois#1{×#1}%
  \def\barre#1{#1}\def\intervalle#1#2{[#1;#2]}\def\intff#1#2{[#1;#2]}%
  \def\textsuperscript#1{#1}\def\dfrac#1#2{#1/#2}\def\frac#1#2{#1/#2}%
  \def\vec#1{vecteur #1}}
\RequirePackage[official]{eurosym}

\tcbset{exercice corrige/.style={
  enhanced,breakable,before upper=\raggedright,
  colback=white,colframe=coulExo,
  boxrule=0.8pt,arc=2pt,
  left=6pt,right=6pt,top=8pt,bottom=5pt,
  before skip=9pt,after skip=4pt,
  coltitle=white,fonttitle=\bfseries\small,
  attach boxed title to top left={xshift=8pt,yshift=-\tcboxedtitleheight/2},
  boxed title style={colback=coulExo,arc=2pt,boxrule=0pt}}}
\newcounter{exo}
\newtcolorbox{exercice}[1][]{exercice corrige,
  phantom={\signet[exercice-\arabic{exo}]{2}{Exercice \arabic{exo}}},
  title={Exercice \theexo\ifx\relax#1\relax\else{} --- #1\fi}}
\newenvironment{exo}[1][\relax]{\refstepcounter{exo}\begin{exercice}[#1]}{\end{exercice}}

\geometry{top=1.8cm,bottom=1cm,left=1cm,right=1cm,headheight=14pt,footskip=0.6cm}
\renewcommand{\lycee@pied}{\thepage}

\renewcommand{\theexo}{\ifnum\value{exo}<10 0\fi\arabic{exo}}
\tcbset{exercice corrige/.style={
  enhanced,breakable,before upper=\raggedright,
  colback=white,colframe=coulExo,
  boxrule=0.9pt,arc=2pt,
  left=8pt,right=8pt,top=8pt,bottom=6pt,
  before skip=8pt,after skip=6pt,
  coltitle=white,fonttitle=\bfseries,
  attach boxed title to top left={xshift=8pt,yshift=-\tcboxedtitleheight/2},
  boxed title style={colback=coulExo,arc=2pt,boxrule=0pt}}}

\newcommand{\entetefiche}[2]{%
  \begin{center}
    \begin{tcolorbox}[enhanced,width=0.92\linewidth,colback=white,colframe=coulPart,
        boxrule=1.4pt,arc=3pt,halign=center,top=6pt,bottom=6pt]
      {\large\bfseries\color{coulPart} CORRIGÉ --- FICHE D'EXERCICES}\\[3pt]
      {\normalsize\bfseries #1}\\[2pt]
      {\normalsize #2}
    \end{tcolorbox}
  \end{center}}

\newcounter{partie}
\newcommand{\partie}[1]{%
  \stepcounter{partie}%
  \besoinplace{8\baselineskip}\vspace{4pt}%
  \begin{tcolorbox}[enhanced,colback=coulPart!8,colframe=coulPart,boxrule=0pt,
      leftrule=3pt,arc=0pt,left=6pt,right=4pt,top=3pt,bottom=3pt,
      before skip=6pt,after skip=2pt]
    \bfseries\color{coulPart}\Roman{partie} --- #1\signet{1}{\Roman{partie} --- #1}
  \end{tcolorbox}}

\setlist[questions,1]{label=\arabic*\textdegree),leftmargin=*,itemsep=3pt,topsep=3pt}
\setlist[questions,2]{label=\alph*),leftmargin=*,itemsep=2pt,topsep=2pt}
\setlist[questions,3]{label=\roman*),leftmargin=*,itemsep=1pt,topsep=1pt}
\newlist{expr}{enumerate}{1}
\setlist[expr]{label={\Alph*\ $=$},leftmargin=*,itemsep=2pt,topsep=3pt}

\newcommand{\res}[1]{{\color{coulRep}#1}}
\newcommand{\rem}[1]{\par\smallskip{\small\itshape\color{black!60}$\rightarrow$ #1}}
\newenvironment{calculs}{\par\smallskip\noindent$\displaystyle\begin{aligned}[t]}%
  {\end{aligned}$\par\smallskip}
\newcommand{\justif}[1]{\quad\text{\small\itshape\color{black!60}#1}}

\setlength{\parskip}{2pt}

\makeatother
\geometry{top=1.8cm,bottom=1.6cm,left=1.8cm,right=1.8cm,headheight=14pt,footskip=30pt}
\entete{Chapitre 4 --- Dérivation --- corrigé}{Première spécialité}
\usepackage{tkz-tab}

\begin{document}

\entetefiche{Chapitre 4 --- Dérivation}{Première spécialité mathématiques}

\partie{Limite en zéro d'une fonction}

\begin{exo}[Conjecturer avec un tableau de valeurs]
\begin{questions}[label=\arabic*.]
  \item $f(-0{,}1)=-3{,}1$ ; $f(-0{,}01)=-3{,}01$ ; $f(0{,}01)=-2{,}99$ ;
        $f(0{,}1)=-2{,}9$.
  \item $f(x)$ semble s'approcher de $\res{-3}$.
  \item Pour $x\ne 0$ :
\begin{calculs}
f(x) &= \frac{\cancel{x}(x-3)}{\cancel{x}} &&\justif{on factorise par $x$}\\
f(x) &= x-3
\end{calculs}
        Quand $x$ tend vers $0$, $x-3$ tend vers $-3$ : $\res{\lim\limits_{x\to 0}f(x)=-3}$.
\end{questions}
\end{exo}

\begin{exo}[Calculer des limites en zéro]
On factorise par $x$ et on simplifie, puis on remplace $x$ par $0$.
\begin{questions}[label=\alph*)]
  \item $f(x)=x+4$ : limite $\res{4}$.
  \item $g(x)=2x-6$ : limite $\res{-6}$.
  \item $h(x)=4x^{2}+1$ : limite $\res{1}$.
  \item $k(x)=3x^{2}-x+8$ : limite $\res{8}$.
\end{questions}
\end{exo}

\begin{exo}[Avec une identité remarquable]
On développe le carré : la constante disparaît, puis on simplifie par $x$.
\begin{questions}[label=\alph*)]
  \item $f(x)=\dfrac{1+2x+x^{2}-1}{x}=2+x$ : limite $\res{2}$.
  \item $g(x)=\dfrac{25+10x+x^{2}-25}{x}=10+x$ : limite $\res{10}$.
  \item $h(x)=\dfrac{x^{2}-4x+4-4}{x}=x-4$ : limite $\res{-4}$.
  \item $k(x)=\dfrac{9-6x+x^{2}-9}{x}=x-6$ : limite $\res{-6}$.
\end{questions}
\end{exo}

\begin{exo}[Limites quand $h$ tend vers $0$]
\begin{questions}[label=\alph*)]
  \item $\dfrac{h^{2}+6h}{h}=h+6$ : limite $\res{6}$.
  \item $\dfrac{4+4h+h^{2}-4}{h}=4+h$ : limite $\res{4}$.
  \item $\dfrac{3+6h+3h^{2}-3}{h}=6+3h$ : limite $\res{6}$.
\end{questions}
\end{exo}

\partie{Nombre dérivé}

\begin{exo}[Nombre dérivé par la définition]
\begin{questions}[label=\alph*)]
  \item
\begin{calculs}
T(h) &= \frac{(3+h)^{2}-9}{h} &&\\
T(h) &= \frac{9+6h+h^{2}-9}{h} &&\\
T(h) &= 6+h &&\\
f'(3) &= \lim_{h\to 0}\,(6+h) &&\\
f'(3) &= \res{6}
\end{calculs}
  \item $g(1)=3$ et $g(1+h)=1+2h+h^{2}+2+2h=3+4h+h^{2}$.
\begin{calculs}
T(h) &= \frac{3+4h+h^{2}-3}{h} &&\\
T(h) &= 4+h &&\\
g'(1) &= \lim_{h\to 0}\,(4+h) &&\\
g'(1) &= \res{4}
\end{calculs}
  \item $k(1)=4$ et $k(1+h)=-(1+2h+h^{2})+5+5h=4+3h-h^{2}$.
\begin{calculs}
T(h) &= \frac{3h-h^{2}}{h} &&\\
T(h) &= 3-h &&\\
k'(1) &= \lim_{h\to 0}\,(3-h) &&\\
k'(1) &= \res{3}
\end{calculs}
  \item $m(-1)=-1$ et $m(-1+h)=2(1-2h+h^{2})-3=-1-4h+2h^{2}$.
\begin{calculs}
T(h) &= \frac{-4h+2h^{2}}{h} &&\\
T(h) &= -4+2h &&\\
m'(-1) &= \lim_{h\to 0}\,(-4+2h) &&\\
m'(-1) &= \res{-4}
\end{calculs}
\end{questions}
\rem{Erreur fréquente : oublier les parenthèses dans $-(1+h)^{2}$.}
\end{exo}

\begin{exo}[Le taux de variation est donné]
On remplace $h$ par $0$ dans $T(h)$.
\begin{questions}[label=\alph*)]
  \item $f'(1)=\lim\limits_{h\to 0}\,(3h+7)=\res{7}$.
  \item $g'(4)=\lim\limits_{h\to 0}\,(h^{2}-h-2)=\res{-2}$.
  \item On divise d'abord par $h$ : $T(h)=\dfrac{5h-2h^{2}}{h}=5-2h$, donc
        $k'(-1)=\res{5}$.
\end{questions}
\end{exo}

\begin{exo}[Avec des fractions]
\begin{questions}[label=\arabic*.]
  \item
\begin{calculs}
f(4+h)-f(4) &= \frac{1}{4+h}-\frac{1}{4} &&\\
f(4+h)-f(4) &= \frac{1\fois{4}}{(4+h)\fois{4}}-\frac{1\fois{(4+h)}}{4\fois{(4+h)}} &&\justif{même dénominateur}\\
f(4+h)-f(4) &= \frac{-h}{4(4+h)}
\end{calculs}
        On divise par $h$, puis on passe à la limite :
\begin{calculs}
T(h) &= \frac{-1}{4(4+h)} &&\\
f'(4) &= \lim_{h\to 0}\frac{-1}{4(4+h)} &&\\
f'(4) &= \res{-\frac{1}{16}}
\end{calculs}
  \item $g(1)=\frac{1}{2}$ et $g(1+h)=\frac{1}{2+h}$.
\begin{calculs}
g(1+h)-g(1) &= \frac{1}{2+h}-\frac{1}{2} &&\\
g(1+h)-g(1) &= \frac{2-(2+h)}{2(2+h)} &&\\
g(1+h)-g(1) &= \frac{-h}{2(2+h)}
\end{calculs}
        On divise par $h$, puis on passe à la limite :
\begin{calculs}
T(h) &= \frac{-1}{2(2+h)} &&\\
g'(1) &= \lim_{h\to 0}\frac{-1}{2(2+h)} &&\\
g'(1) &= \res{-\frac{1}{4}}
\end{calculs}
\end{questions}
\end{exo}

\begin{exo}[Pour un nombre $a$ quelconque]
\begin{questions}[label=\arabic*.]
  \item
\begin{calculs}
f(a+h)-f(a) &= (a+h)^{2}-4(a+h)-a^{2}+4a &&\\
f(a+h)-f(a) &= a^{2}+2ah+h^{2}-4a-4h-a^{2}+4a &&\\
f(a+h)-f(a) &= 2ah+h^{2}-4h &&\\
T(h) &= 2a+h-4 &&\justif{on divise par $h$}
\end{calculs}
  \item $f'(a)=\lim\limits_{h\to 0}\,(2a-4+h)=\res{2a-4}$.
  \item $f'(0)=\res{-4}$ ; $f'(2)=\res{0}$ ; $f'(5)=\res{6}$.
\end{questions}
\end{exo}

\partie{Tangente à une courbe}

\begin{exo}[Lecture graphique]
\begin{questions}[label=\arabic*.]
  \item $f(-1)=\res{2}$ ; $f(1)=\res{0}$ ; $f(3)=\res{-2}$.
  \item Les tangentes en $A$ et $C$ sont horizontales : $f'(-1)=\res{0}$ et
        $f'(3)=\res{0}$. La tangente en $B$ passe par $B(1\,;0)$ et par $(3\,;-3)$ :
        quand on avance de $2$, elle descend de $3$, donc $f'(1)=\res{-\frac{3}{2}}$.
  \item
\begin{calculs}
y &= f'(1)(x-1)+f(1) &&\\
y &= -\frac{3}{2}(x-1)+0 &&\\
y &= \res{-\frac{3}{2}x+\frac{3}{2}}
\end{calculs}
\end{questions}
\end{exo}

\begin{exo}[Équation de la tangente]
On applique $y=f'(a)(x-a)+f(a)$.
\begin{questions}[label=\alph*)]
  \item $y=4(x-2)+3$, soit $\res{y=4x-5}$.
  \item $y=-3(x+1)+2$, soit $\res{y=-3x-1}$.
  \item $y=\frac{1}{2}(x-0)+5$, soit $\res{y=\frac{1}{2}x+5}$.
  \item $y=0\times(x-4)-1$, soit $\res{y=-1}$ (tangente horizontale).
\end{questions}
\end{exo}

\begin{exo}[Nombre dérivé, puis tangente]
\begin{questions}[label=\arabic*.]
  \item $f(1)=0$ et $f(1+h)=1+2h+h^{2}+1+h-2=3h+h^{2}$.
\begin{calculs}
T(h) &= \frac{3h+h^{2}}{h} &&\\
T(h) &= 3+h &&\\
f'(1) &= \lim_{h\to 0}\,(3+h) &&\\
f'(1) &= 3 &&\\
y &= 3(x-1)+0 &&\justif{$y=f'(1)(x-1)+f(1)$}\\
y &= \res{3x-3}
\end{calculs}
  \item $g(-1)=-1$.
\begin{calculs}
g(-1+h)-g(-1) &= \frac{1}{-1+h}+1 &&\\
g(-1+h)-g(-1) &= \frac{1+(-1+h)}{-1+h} &&\justif{même dénominateur}\\
g(-1+h)-g(-1) &= \frac{h}{-1+h} &&\\
T(h) &= \frac{1}{-1+h} &&\\
g'(-1) &= \lim_{h\to 0}\frac{1}{-1+h} &&\\
g'(-1) &= -1 &&\\
y &= -(x+1)-1 &&\\
y &= \res{-x-2}
\end{calculs}
\end{questions}
\end{exo}

\begin{exo}[Tangente passant par un point]
$f'(1)$ est le coefficient directeur de la tangente, donc de la droite $(AB)$.
\begin{questions}[label=\arabic*.]
  \item
\begin{calculs}
f'(1) &= \frac{4-(-2)}{3-1} &&\justif{$\frac{y_{B}-y_{A}}{x_{B}-x_{A}}$}\\
f'(1) &= \res{3} &&\\
y &= 3(x-1)-2 &&\\
y &= \res{3x-5}
\end{calculs}
  \item
\begin{calculs}
g'(-2) &= \frac{-3-1}{0-(-2)} &&\\
g'(-2) &= \res{-2} &&\\
y &= -2(x+2)+1 &&\\
y &= \res{-2x-3}
\end{calculs}
\end{questions}
\end{exo}

\partie{Dérivées des fonctions polynômes}

\begin{exo}[Dériver un polynôme]
\begin{questions}[label=\alph*)]
  \item $f'(x)=\res{0}$ (constante)
  \item $g'(x)=\res{5}$
  \item $h'(x)=\res{6x^{5}}$
  \item $k'(x)=\res{6x-4}$
  \item $m'(x)=\res{-6x^{2}+2x-5}$
  \item $p'(x)=\res{4x^{3}-6x^{2}+7}$
\end{questions}
\rem{Le coefficient se garde : la dérivée de $-2x^{3}$ est $-2\times 3x^{2}=-6x^{2}$.}
\end{exo}

\begin{exo}[Avec des fractions]
\begin{questions}[label=\alph*)]
  \item $f'(x)=\frac{1}{2}\times 2x-3=\res{x-3}$
  \item $g'(x)=\frac{1}{3}\times 3x^{2}-2x=\res{x^{2}-2x}$
  \item $h'(x)=\frac{1}{4}\times 4x^{3}+\frac{2}{3}\times 3x^{2}=\res{x^{3}+2x^{2}}$
  \item $k(x)=\frac{1}{3}x^{3}-2x$, donc $k'(x)=\res{x^{2}-2}$
\end{questions}
\end{exo}

\begin{exo}[Développer d'abord]
\begin{questions}[label=\alph*)]
  \item $f(x)=x^{2}+8x+16$, donc $f'(x)=\res{2x+8}$
  \item $g(x)=9x^{2}-6x+1$, donc $g'(x)=\res{18x-6}$
  \item $h(x)=x^{3}-5x$, donc $h'(x)=\res{3x^{2}-5}$
  \item $k(x)=x^{2}-4$, donc $k'(x)=\res{2x}$
\end{questions}
\end{exo}

\begin{exo}[Entraînement : polynômes]
\begin{questions}[label=\alph*)]
  \item $f'(x)=\res{3x^{2}-14x+2}$
  \item $g'(x)=\res{-20x^{4}+6x}$
  \item $h'(x)=\res{24x^{3}-3x^{2}+1}$
  \item $k'(x)=\frac{2}{5}\times 5x^{4}-1=\res{2x^{4}-1}$
  \item $m'(x)=0{,}5\times 4x^{3}-1{,}5\times 2x=\res{2x^{3}-3x}$
  \item $p(x)=\frac{1}{4}x^{2}-\frac{3}{4}$, donc $p'(x)=\res{\frac{1}{2}x}$
  \item $q'(x)=\res{7x^{6}-5x^{4}+3x^{2}}$
  \item $r(x)=3x^{2}-6x+15$, donc $r'(x)=\res{6x-6}$
\end{questions}
\end{exo}

\begin{exo}[Fonction dérivée et nombres dérivés]
\begin{questions}[label=\arabic*.]
  \item $f'(x)=\res{3x^{2}-6x}$.
  \item $f'(0)=\res{0}$ ; $f'(1)=3-6=\res{-3}$ ; $f'(-1)=3+6=\res{9}$ ;
        $f'(2)=12-12=\res{0}$.
  \item $f(1)=1-3+2=0$.
\begin{calculs}
y &= f'(1)(x-1)+f(1) &&\\
y &= -3(x-1)+0 &&\\
y &= \res{-3x+3}
\end{calculs}
  \item La tangente est horizontale quand $f'(x)=0$ : $3x(x-2)=0$, soit $x=0$ ou
        $x=2$. Avec $f(0)=2$ et $f(2)=8-12+2=-2$ : aux points $\res{(0\,;2)}$ et
        $\res{(2\,;-2)}$.
\end{questions}
\end{exo}

\partie{Dérivées des autres fonctions de référence}

\begin{exo}[Nombres dérivés des fonctions de référence]
\begin{questions}[label=\alph*)]
  \item $f'(x)=-\dfrac{1}{x^{2}}$ : $f'(5)=\res{-\frac{1}{25}}$ et
        $f'\bigl(\frac{1}{2}\bigr)=-\dfrac{1}{\frac{1}{4}}=\res{-4}$.
  \item $g'(x)=\dfrac{1}{2\sqrt{x}}$ : $g'(16)=\dfrac{1}{2\times 4}=\res{\frac{1}{8}}$
        et $g'\bigl(\frac{1}{4}\bigr)=\dfrac{1}{2\times\frac{1}{2}}=\res{1}$.
  \item $k'(x)=-\dfrac{2}{x^{3}}$ : $k'(1)=\res{-2}$ et
        $k'(-2)=-\dfrac{2}{-8}=\res{\frac{1}{4}}$.
\end{questions}
\end{exo}

\begin{exo}[Dériver avec les fonctions de référence]
\begin{questions}[label=\alph*)]
  \item $f'(x)=4\times\left(-\dfrac{1}{x^{2}}\right)=\res{-\dfrac{4}{x^{2}}}$
  \item $g'(x)=-3\times\dfrac{1}{2\sqrt{x}}=\res{-\dfrac{3}{2\sqrt{x}}}$
  \item $h'(x)=\res{-\dfrac{4}{x^{5}}}$
  \item $k'(x)=5\times\left(-\dfrac{3}{x^{4}}\right)=\res{-\dfrac{15}{x^{4}}}$
  \item $m'(x)=\res{-2\sin x}$
  \item $p'(x)=\res{-\cos x}$
\end{questions}
\end{exo}

\begin{exo}[Entraînement : fonctions de référence]
\begin{questions}[label=\alph*)]
  \item $f'(x)=-7\times\left(-\dfrac{1}{x^{2}}\right)=\res{\dfrac{7}{x^{2}}}$
  \item $g'(x)=8\times\dfrac{1}{2\sqrt{x}}=\res{\dfrac{4}{\sqrt{x}}}$
  \item $h'(x)=2\times\left(-\dfrac{2}{x^{3}}\right)=\res{-\dfrac{4}{x^{3}}}$
  \item $k'(x)=-1\times\left(-\dfrac{5}{x^{6}}\right)=\res{\dfrac{5}{x^{6}}}$
  \item $m'(x)=\dfrac{1}{3}\times\dfrac{1}{2\sqrt{x}}=\res{\dfrac{1}{6\sqrt{x}}}$
  \item $p'(x)=-4\times(-\sin x)=\res{4\sin x}$
  \item $q(x)=\dfrac{1}{2}\times\dfrac{1}{x}$, donc
        $q'(x)=\dfrac{1}{2}\times\left(-\dfrac{1}{x^{2}}\right)=\res{-\dfrac{1}{2x^{2}}}$
  \item $r'(x)=\res{\dfrac{1}{2}\cos x}$
\end{questions}
\rem{Deux signes « $-$ » donnent un « $+$ » : c'est l'erreur la plus fréquente en a), d)
et f).}
\end{exo}

\begin{exo}[Tangentes aux courbes de référence]
\begin{questions}[label=\arabic*.]
  \item $f(1)=1$ et $f'(1)=-1$.
\begin{calculs}
y &= -1\times(x-1)+1 &&\\
y &= \res{-x+2}
\end{calculs}
  \item $g(4)=2$ et $g'(4)=\dfrac{1}{2\times 2}=\dfrac{1}{4}$.
\begin{calculs}
y &= \frac{1}{4}(x-4)+2 &&\\
y &= \res{\frac{1}{4}x+1}
\end{calculs}
  \item $f'(x)=-\dfrac{1}{x^{2}}$ est toujours strictement négatif (un carré non nul est
        positif) : $f'(x)$ ne vaut jamais $0$, donc \res{aucune tangente n'est
        horizontale}.
\end{questions}
\end{exo}

\partie{Dérivée d'une somme et d'un produit}

\begin{exo}[Dériver une somme]
On dérive chaque terme.
\begin{questions}[label=\alph*)]
  \item $f'(x)=\res{3x^{2}-\dfrac{1}{x^{2}}}$
  \item $g'(x)=\res{4x-\dfrac{1}{2\sqrt{x}}}$
  \item $h'(x)=\res{-\dfrac{3}{x^{2}}+5}$
  \item $k'(x)=\res{1-\sin x}$
\end{questions}
\end{exo}

\begin{exo}[Dériver un produit]
\begin{questions}[label=\alph*)]
  \item $u(x)=3x-2$, $u'(x)=3$ ; $v(x)=x+4$, $v'(x)=1$.
\begin{calculs}
f'(x) &= 3(x+4)+(3x-2)\times 1 &&\justif{$u'v+uv'$}\\
f'(x) &= 3x+12+3x-2 &&\\
f'(x) &= \res{6x+10}
\end{calculs}
  \item $u(x)=x^{2}+1$, $u'(x)=2x$ ; $v(x)=2x-3$, $v'(x)=2$.
\begin{calculs}
g'(x) &= 2x(2x-3)+(x^{2}+1)\times 2 &&\\
g'(x) &= 4x^{2}-6x+2x^{2}+2 &&\\
g'(x) &= \res{6x^{2}-6x+2}
\end{calculs}
  \item $u(x)=x+1$, $u'(x)=1$ ; $v(x)=\sqrt{x}$, $v'(x)=\dfrac{1}{2\sqrt{x}}$.
\begin{calculs}
h'(x) &= 1\times\sqrt{x}+(x+1)\times\frac{1}{2\sqrt{x}} &&\\
h'(x) &= \frac{2x}{2\sqrt{x}}+\frac{x+1}{2\sqrt{x}} &&\justif{$\sqrt{x}=\frac{\sqrt{x}\times 2\sqrt{x}}{2\sqrt{x}}=\frac{2x}{2\sqrt{x}}$}\\
h'(x) &= \res{\frac{3x+1}{2\sqrt{x}}}
\end{calculs}
  \item $u(x)=x^{2}$, $u'(x)=2x$ ; $v(x)=\cos x$, $v'(x)=-\sin x$.
\begin{calculs}
k'(x) &= 2x\cos x+x^{2}\times(-\sin x) &&\\
k'(x) &= \res{2x\cos x-x^{2}\sin x}
\end{calculs}
\end{questions}
\end{exo}

\begin{exo}[Entraînement : sommes et produits]
\begin{questions}[label=\alph*)]
  \item $f'(x)=2x-3+4\times\left(-\dfrac{1}{x^{2}}\right)=\res{2x-3-\dfrac{4}{x^{2}}}$
  \item $g'(x)=4\times\dfrac{1}{2\sqrt{x}}-6x^{2}=\res{\dfrac{2}{\sqrt{x}}-6x^{2}}$
  \item $u(x)=x^{2}-4$, $u'(x)=2x$ ; $v(x)=3x+1$, $v'(x)=3$.
\begin{calculs}
h'(x) &= 2x(3x+1)+(x^{2}-4)\times 3 &&\\
h'(x) &= 6x^{2}+2x+3x^{2}-12 &&\\
h'(x) &= \res{9x^{2}+2x-12}
\end{calculs}
  \item $u(x)=2x-1$, $u'(x)=2$ ; $v(x)=\sqrt{x}$, $v'(x)=\dfrac{1}{2\sqrt{x}}$.
\begin{calculs}
k'(x) &= 2\sqrt{x}+(2x-1)\times\frac{1}{2\sqrt{x}} &&\\
k'(x) &= \frac{4x}{2\sqrt{x}}+\frac{2x-1}{2\sqrt{x}} &&\justif{$2\sqrt{x}=\frac{4x}{2\sqrt{x}}$}\\
k'(x) &= \res{\frac{6x-1}{2\sqrt{x}}}
\end{calculs}
  \item $u(x)=x+2$, $u'(x)=1$ ; $v(x)=\cos x$, $v'(x)=-\sin x$.
\begin{calculs}
m'(x) &= 1\times\cos x+(x+2)(-\sin x) &&\\
m'(x) &= \res{\cos x-(x+2)\sin x}
\end{calculs}
  \item $u(x)=\dfrac{1}{x}$, $u'(x)=-\dfrac{1}{x^{2}}$ ; $v(x)=x+5$, $v'(x)=1$.
\begin{calculs}
p'(x) &= -\frac{1}{x^{2}}(x+5)+\frac{1}{x}\times 1 &&\\
p'(x) &= \frac{-x-5}{x^{2}}+\frac{x}{x^{2}} &&\justif{même dénominateur}\\
p'(x) &= \res{-\frac{5}{x^{2}}}
\end{calculs}
        \rem{Plus court : $p(x)=1+\dfrac{5}{x}$, puis on dérive la somme.}
  \item $u(x)=x^{2}+x$, $u'(x)=2x+1$ ; $v(x)=x^{2}-x$, $v'(x)=2x-1$.
\begin{calculs}
q'(x) &= (2x+1)(x^{2}-x)+(x^{2}+x)(2x-1) &&\\
q'(x) &= 2x^{3}-x^{2}-x+2x^{3}+x^{2}-x &&\\
q'(x) &= \res{4x^{3}-2x}
\end{calculs}
        \rem{Vérification : $q(x)=x^{4}-x^{2}$, de dérivée $4x^{3}-2x$.}
  \item $u(x)=\sin x$, $u'(x)=\cos x$ ; $v(x)=\cos x$, $v'(x)=-\sin x$.
\begin{calculs}
r'(x) &= \cos x\times\cos x+\sin x\times(-\sin x) &&\\
r'(x) &= \res{(\cos x)^{2}-(\sin x)^{2}}
\end{calculs}
\end{questions}
\end{exo}

\begin{exo}[Deux façons de dériver]
\begin{questions}[label=\arabic*.]
  \item $u(x)=2x+5$, $u'(x)=2$ ; $v(x)=1-3x$, $v'(x)=-3$.
\begin{calculs}
f'(x) &= 2(1-3x)+(2x+5)(-3) &&\\
f'(x) &= 2-6x-6x-15 &&\\
f'(x) &= \res{-12x-13}
\end{calculs}
  \item
\begin{calculs}
f(x) &= 2x-6x^{2}+5-15x &&\\
f(x) &= -6x^{2}-13x+5 &&\\
f'(x) &= \res{-12x-13}
\end{calculs}
  \item On trouve bien la même dérivée.
\end{questions}
\end{exo}

\begin{exo}[Produit et tangente]
\begin{questions}[label=\arabic*.]
  \item $u(x)=x-3$, $u'(x)=1$ ; $v(x)=\sqrt{x}$, $v'(x)=\dfrac{1}{2\sqrt{x}}$.
\begin{calculs}
f'(x) &= \sqrt{x}+\frac{x-3}{2\sqrt{x}} &&\\
f'(x) &= \frac{2x}{2\sqrt{x}}+\frac{x-3}{2\sqrt{x}} &&\\
f'(x) &= \frac{3x-3}{2\sqrt{x}}
\end{calculs}
  \item $f(4)=1\times 2=\res{2}$ et $f'(4)=\dfrac{12-3}{2\times 2}=\res{\dfrac{9}{4}}$.
  \item
\begin{calculs}
y &= \frac{9}{4}(x-4)+2 &&\\
y &= \frac{9}{4}x-9+2 &&\\
y &= \res{\frac{9}{4}x-7}
\end{calculs}
\end{questions}
\end{exo}

\partie{Dérivée d'un inverse et d'un quotient}

\begin{exo}[Dériver un inverse]
On applique $\left(\dfrac{1}{v}\right)'=-\dfrac{v'}{v^{2}}$.
\begin{questions}[label=\alph*)]
  \item $v(x)=2x+3$, $v'(x)=2$ : $f'(x)=\res{-\dfrac{2}{(2x+3)^{2}}}$
  \item $v(x)=x^{2}+4$, $v'(x)=2x$ : $g'(x)=\res{-\dfrac{2x}{(x^{2}+4)^{2}}}$
  \item $h(x)=3\times\dfrac{1}{x-1}$ : $h'(x)=3\times\left(-\dfrac{1}{(x-1)^{2}}\right)
        =\res{-\dfrac{3}{(x-1)^{2}}}$
  \item $v(x)=x^{2}-2x$, $v'(x)=2x-2$ : $k'(x)=\res{-\dfrac{2x-2}{(x^{2}-2x)^{2}}}$
\end{questions}
\end{exo}

\begin{exo}[Dériver un quotient]
On applique $\left(\dfrac{u}{v}\right)'=\dfrac{u'v-uv'}{v^{2}}$.
\begin{questions}[label=\alph*)]
  \item $u(x)=x+3$, $u'(x)=1$ ; $v(x)=x-2$, $v'(x)=1$.
\begin{calculs}
f'(x) &= \frac{1\times(x-2)-(x+3)\times 1}{(x-2)^{2}} &&\\
f'(x) &= \frac{x-2-x-3}{(x-2)^{2}} &&\\
f'(x) &= \res{\frac{-5}{(x-2)^{2}}}
\end{calculs}
  \item $u(x)=4x-1$, $u'(x)=4$ ; $v(x)=2x+1$, $v'(x)=2$.
\begin{calculs}
g'(x) &= \frac{4(2x+1)-(4x-1)\times 2}{(2x+1)^{2}} &&\\
g'(x) &= \frac{8x+4-8x+2}{(2x+1)^{2}} &&\\
g'(x) &= \res{\frac{6}{(2x+1)^{2}}}
\end{calculs}
  \item $u(x)=x^{2}$, $u'(x)=2x$ ; $v(x)=x-3$, $v'(x)=1$.
\begin{calculs}
h'(x) &= \frac{2x(x-3)-x^{2}\times 1}{(x-3)^{2}} &&\\
h'(x) &= \res{\frac{x^{2}-6x}{(x-3)^{2}}}
\end{calculs}
  \item $u(x)=x^{2}-1$, $u'(x)=2x$ ; $v(x)=x^{2}+1$, $v'(x)=2x$.
\begin{calculs}
k'(x) &= \frac{2x(x^{2}+1)-(x^{2}-1)\times 2x}{(x^{2}+1)^{2}} &&\\
k'(x) &= \frac{2x^{3}+2x-2x^{3}+2x}{(x^{2}+1)^{2}} &&\\
k'(x) &= \res{\frac{4x}{(x^{2}+1)^{2}}}
\end{calculs}
\end{questions}
\rem{Le « $-$ » devant $uv'$ change tous les signes : on met $uv'$ entre parenthèses.}
\end{exo}

\begin{exo}[Entraînement : inverses et quotients]
\begin{questions}[label=\alph*)]
  \item $v(x)=3-x$, $v'(x)=-1$ : $f'(x)=-\dfrac{-1}{(3-x)^{2}}=\res{\dfrac{1}{(3-x)^{2}}}$
  \item $g(x)=-2\times\dfrac{1}{x^{2}+x+1}$ et $v'(x)=2x+1$ :
        $g'(x)=-2\times\left(-\dfrac{2x+1}{(x^{2}+x+1)^{2}}\right)
        =\res{\dfrac{4x+2}{(x^{2}+x+1)^{2}}}$
  \item $u'(x)=5$ ; $v'(x)=1$.
\begin{calculs}
h'(x) &= \frac{5(x-1)-(5x+2)}{(x-1)^{2}} &&\\
h'(x) &= \res{\frac{-7}{(x-1)^{2}}}
\end{calculs}
  \item $u'(x)=-2$ ; $v'(x)=1$.
\begin{calculs}
k'(x) &= \frac{-2(x+3)-(1-2x)}{(x+3)^{2}} &&\\
k'(x) &= \frac{-2x-6-1+2x}{(x+3)^{2}} &&\\
k'(x) &= \res{\frac{-7}{(x+3)^{2}}}
\end{calculs}
  \item $u(x)=x^{2}+1$, $u'(x)=2x$ ; $v(x)=x$, $v'(x)=1$.
\begin{calculs}
m'(x) &= \frac{2x\times x-(x^{2}+1)}{x^{2}} &&\\
m'(x) &= \res{\frac{x^{2}-1}{x^{2}}}
\end{calculs}
  \item $u(x)=3x$, $u'(x)=3$ ; $v(x)=x^{2}+1$, $v'(x)=2x$.
\begin{calculs}
p'(x) &= \frac{3(x^{2}+1)-3x\times 2x}{(x^{2}+1)^{2}} &&\\
p'(x) &= \res{\frac{3-3x^{2}}{(x^{2}+1)^{2}}}
\end{calculs}
  \item $u(x)=x^{2}-4x$, $u'(x)=2x-4$ ; $v(x)=x+2$, $v'(x)=1$.
\begin{calculs}
q'(x) &= \frac{(2x-4)(x+2)-(x^{2}-4x)}{(x+2)^{2}} &&\\
q'(x) &= \frac{2x^{2}-8-x^{2}+4x}{(x+2)^{2}} &&\justif{$(2x-4)(x+2)=2x^{2}-8$}\\
q'(x) &= \res{\frac{x^{2}+4x-8}{(x+2)^{2}}}
\end{calculs}
  \item $u(x)=\sqrt{x}$, $u'(x)=\dfrac{1}{2\sqrt{x}}$ ; $v(x)=x+1$, $v'(x)=1$.
\begin{calculs}
r'(x) &= \frac{\frac{1}{2\sqrt{x}}(x+1)-\sqrt{x}}{(x+1)^{2}} &&\\
r'(x) &= \frac{(x+1)-2x}{2\sqrt{x}\,(x+1)^{2}} &&\justif{on multiplie en haut et en bas par $2\sqrt{x}$}\\
r'(x) &= \res{\frac{1-x}{2\sqrt{x}\,(x+1)^{2}}}
\end{calculs}
\end{questions}
\end{exo}

\begin{exo}[Choisir la bonne formule]
\begin{questions}[label=\alph*)]
  \item Somme : $f'(x)=12x^{3}-2\times\left(-\dfrac{1}{x^{2}}\right)
        =\res{12x^{3}+\dfrac{2}{x^{2}}}$
  \item Produit, avec $u'(x)=1$ et $v'(x)=2x$.
\begin{calculs}
g'(x) &= 1\times(x^{2}+2)+(x-1)\times 2x &&\\
g'(x) &= x^{2}+2+2x^{2}-2x &&\\
g'(x) &= \res{3x^{2}-2x+2}
\end{calculs}
  \item Inverse (multiplié par $4$), avec $v'(x)=2x$ :
        $h'(x)=4\times\left(-\dfrac{2x}{(x^{2}+3)^{2}}\right)=\res{-\dfrac{8x}{(x^{2}+3)^{2}}}$
  \item Quotient, avec $u'(x)=1$ et $v'(x)=2$.
\begin{calculs}
k'(x) &= \frac{1\times(2x+1)-(x-4)\times 2}{(2x+1)^{2}} &&\\
k'(x) &= \frac{2x+1-2x+8}{(2x+1)^{2}} &&\\
k'(x) &= \res{\frac{9}{(2x+1)^{2}}}
\end{calculs}
  \item Produit, avec $u(x)=x^{2}$ et $v(x)=\sqrt{x}$.
\begin{calculs}
m'(x) &= 2x\sqrt{x}+x^{2}\times\frac{1}{2\sqrt{x}} &&\\
m'(x) &= 2x\sqrt{x}+\frac{x\sqrt{x}}{2} &&\justif{$\frac{x^{2}}{\sqrt{x}}=x\sqrt{x}$}\\
m'(x) &= \res{\frac{5}{2}x\sqrt{x}}
\end{calculs}
  \item Somme : $p'(x)=\res{\dfrac{1}{2\sqrt{x}}+\cos x}$
\end{questions}
\end{exo}

\begin{exo}[Tangentes à une hyperbole]
\begin{questions}[label=\arabic*.]
  \item $u(x)=2x-4$, $u'(x)=2$ ; $v(x)=x+1$, $v'(x)=1$.
\begin{calculs}
f'(x) &= \frac{2(x+1)-(2x-4)\times 1}{(x+1)^{2}} &&\\
f'(x) &= \frac{2x+2-2x+4}{(x+1)^{2}} &&\\
f'(x) &= \frac{6}{(x+1)^{2}}
\end{calculs}
  \item $f(x)=0$ quand $2x-4=0$, soit $x=2$ : $\res{A(2\,;0)}$.
        $f'(2)=\dfrac{6}{9}=\dfrac{2}{3}$.
\begin{calculs}
y &= \frac{2}{3}(x-2)+0 &&\\
y &= \res{\frac{2}{3}x-\frac{4}{3}}
\end{calculs}
  \item $f(0)=-4$ : $\res{B(0\,;-4)}$. $f'(0)=6$, donc la tangente en $B$ a pour équation
        $\res{y=6x-4}$.
\end{questions}
\end{exo}

\begin{exo}[Deux fonctions, une même dérivée]
\begin{questions}[label=\arabic*.]
  \item
\begin{calculs}
f'(x) &= \frac{2(x-1)-(2x+1)\times 1}{(x-1)^{2}} &&\\
f'(x) &= \res{\frac{-3}{(x-1)^{2}}} &&\\
g'(x) &= 3\times\left(-\frac{1}{(x-1)^{2}}\right) &&\\
g'(x) &= \res{\frac{-3}{(x-1)^{2}}}
\end{calculs}
        $f$ et $g$ ont \res{la même dérivée}.
  \item
\begin{calculs}
f(x)-g(x) &= \frac{2x+1-3}{x-1} &&\\
f(x)-g(x) &= \frac{2(x-1)}{x-1} &&\\
f(x)-g(x) &= \res{2}
\end{calculs}
        $f=g+2$ : ajouter une constante ne change pas la dérivée (la dérivée de $2$ est
        $0$). Les deux courbes se déduisent l'une de l'autre par une translation
        verticale : elles ont les mêmes pentes.
\end{questions}
\end{exo}

\partie{Approximation linéaire}

\begin{exo}[Valeurs approchées sans calculatrice]
\begin{questions}[label=\arabic*.]
  \item $f(16)=4$ et $f'(16)=\dfrac{1}{2\times 4}=\dfrac{1}{8}$.
\begin{calculs}
\sqrt{16{,}08} &\approx 4+\frac{1}{8}\times 0{,}08 &&\\
\sqrt{16{,}08} &\approx \res{4{,}01}
\end{calculs}
  \item $f(1)=1$ et $f'(x)=4x^{3}$, donc $f'(1)=4$ :
        $1{,}03^{4}\approx 1+4\times 0{,}03=\res{1{,}12}$.
  \item $f(x)=\dfrac{1}{x}$, $a=1$ et $h=-0{,}02$ ; $f(1)=1$ et $f'(1)=-1$ :
        $\dfrac{1}{0{,}98}\approx 1+(-1)\times(-0{,}02)=\res{1{,}02}$.
  \item $f(x)=x^{3}$, $a=2$ et $h=0{,}01$ ; $f(2)=8$ et $f'(2)=3\times 4=12$ :
        $2{,}01^{3}\approx 8+12\times 0{,}01=\res{8{,}12}$.
  \item La calculatrice donne $4{,}00999\ldots$ ; $1{,}12550\ldots$ ; $1{,}02040\ldots$ ;
        $8{,}120601$ : les valeurs approchées sont justes au centième.
\end{questions}
\rem{Le point clé : choisir $a$ « rond » (où l'on sait calculer $f(a)$ et $f'(a)$) et
$h$ petit.}
\end{exo}

\begin{exo}[Augmentations et baisses successives]
\begin{questions}[label=\arabic*.]
  \item Avec $f(x)=x^{4}$, $a=1$ et $h=0{,}015$ :
        $1{,}015^{4}\approx 1+4\times 0{,}015=\res{1{,}06}$. Hausse d'environ
        $\res{6~\%}$ (la calculatrice donne $1{,}0614$).
  \item Avec $f(x)=x^{5}$, $a=1$ et $h=-0{,}02$ :
        $0{,}98^{5}\approx 1+5\times(-0{,}02)=\res{0{,}90}$. Baisse d'environ
        $\res{10~\%}$ (la calculatrice donne $0{,}9039$).
\end{questions}
\rem{Pour de petits pourcentages, $n$ hausses de $t~\%$ font environ $n\times t~\%$.}
\end{exo}

\begin{exo}[Courbe et tangente]
\begin{questions}[label=\arabic*.]
  \item $f(2)=4+6=\res{10}$ ; $f'(x)=2x+3$, donc $f'(2)=\res{7}$.
  \item $f(2{,}1)\approx 10+7\times 0{,}1=\res{10{,}7}$. Valeur exacte :
        $f(2{,}1)=4{,}41+6{,}3=10{,}71$ ; l'écart est $\res{0{,}01}$.
  \item $\sin 0=0$ et $\sin'(0)=\cos 0=1$ :
        $\sin(0{,}05)\approx 0+1\times 0{,}05=\res{0{,}05}$.
\end{questions}
\end{exo}

\partie{Étude des variations d'une fonction}

\begin{exo}[Du signe de $f'$ aux variations de $f$]
\begin{questions}[label=\arabic*.]
  \item $f'$ négative : $f$ décroissante ; $f'$ positive : $f$ croissante.\\[2pt]
\begin{tikzpicture}[scale=0.8]
\tkzTabInit[lgt=1.6,espcl=2]{$x$/0.8,$f'(x)$/0.8,$f$/1.6}{$-4$,$-1$,$3$,$5$}
\tkzTabLine{,-,z,+,z,-,}
\tkzTabVar{+/$2$,-/$-3$,+/$4$,-/$1$}
\end{tikzpicture}
  \item $f$ est croissante sur $\intff{-1}{3}$ et $0<2$, donc $\res{f(0)\le f(2)}$
        (et même $f(0)<f(2)$, car $f'(x)>0$ entre $-1$ et $3$).
\end{questions}
\end{exo}

\begin{exo}[Variations de polynômes]
\begin{questions}[label=\alph*)]
  \item $f'(x)=2x-4$ s'annule en $2$ ; $f(2)=4-8+7=3$.\\[2pt]
\begin{tikzpicture}[scale=0.8]
\tkzTabInit[lgt=1.6,espcl=2.2]{$x$/0.8,$f'(x)$/0.8,$f$/1.6}{$-\infty$,$2$,$+\infty$}
\tkzTabLine{,-,z,+,}
\tkzTabVar{+/{},-/$3$,+/{}}
\end{tikzpicture}
  \item $g'(x)=-6x+12$ s'annule en $2$ ; $g(2)=-12+24-5=7$.\\[2pt]
\begin{tikzpicture}[scale=0.8]
\tkzTabInit[lgt=1.6,espcl=2.2]{$x$/0.8,$g'(x)$/0.8,$g$/1.6}{$-\infty$,$2$,$+\infty$}
\tkzTabLine{,+,z,-,}
\tkzTabVar{-/{},+/$7$,-/{}}
\end{tikzpicture}
  \item
\begin{calculs}
h'(x) &= 6x^{2}-6x-12 &&\\
\Delta &= (-6)^{2}-4\times 6\times(-12)=324 &&\justif{$\sqrt{324}=18$}\\
x_{1} &= \frac{6-18}{12}=-1 &&\\
x_{2} &= \frac{6+18}{12}=2 &&\\
h(-1) &= -2-3+12+1=8 &&\\
h(2) &= 16-12-24+1=-19
\end{calculs}
        $h'(x)$ est du signe de $a=6$, sauf entre les racines.\\[2pt]
\begin{tikzpicture}[scale=0.8]
\tkzTabInit[lgt=1.6,espcl=2.2]{$x$/0.8,$h'(x)$/0.8,$h$/1.6}{$-\infty$,$-1$,$2$,$+\infty$}
\tkzTabLine{,+,z,-,z,+,}
\tkzTabVar{-/{},+/$8$,-/$-19$,+/{}}
\end{tikzpicture}
  \item $k'(x)=3x^{2}+3$ : un carré est positif, donc $k'(x)\ge 3>0$. $k$ est
        \res{strictement croissante sur $\R$}.
\end{questions}
\end{exo}

\begin{exo}[Variations d'un quotient]
\begin{questions}[label=\arabic*.]
  \item $u(x)=x^{2}+3$, $u'(x)=2x$ ; $v(x)=x+1$, $v'(x)=1$.
\begin{calculs}
f'(x) &= \frac{2x(x+1)-(x^{2}+3)\times 1}{(x+1)^{2}} &&\\
f'(x) &= \frac{2x^{2}+2x-x^{2}-3}{(x+1)^{2}} &&\\
f'(x) &= \frac{x^{2}+2x-3}{(x+1)^{2}}
\end{calculs}
  \item $(x+1)^{2}>0$ : $f'(x)$ a le signe de $x^{2}+2x-3$, de racines $-3$ et $1$
        ($\Delta=16$). Sur $\intoo{-1}{+\infty}$, seule la racine $1$ compte :
        $\res{f'(x)<0}$ sur $\intoo{-1}{1}$ et $\res{f'(x)>0}$ sur $\intoo{1}{+\infty}$.
  \item $f(1)=\dfrac{4}{2}=2$.\\[2pt]
\begin{tikzpicture}[scale=0.8]
\tkzTabInit[lgt=1.6,espcl=2.2]{$x$/0.8,$f'(x)$/0.8,$f$/1.6}{$-1$,$1$,$+\infty$}
\tkzTabLine{d,-,z,+,}
\tkzTabVar{D+/{},-/$2$,+/{}}
\end{tikzpicture}
\end{questions}
\end{exo}

\begin{exo}[Lire le signe de la dérivée]
\begin{questions}[label=\arabic*.]
  \item La courbe descend, monte, descend puis monte ; $g'$ s'annule aux tangentes
        horizontales.\\[2pt]
\begin{tikzpicture}[scale=0.8]
\tkzTabInit[lgt=1.6,espcl=1.8]{$x$/0.8,$g'(x)$/0.8}{$-4$,$-2$,$1$,$3$,$4$}
\tkzTabLine{,-,z,+,z,-,z,+,}
\end{tikzpicture}
  \item
\begin{tikzpicture}[scale=0.8]
\tkzTabInit[lgt=1.6,espcl=1.8]{$x$/0.8,$g$/1.6}{$-4$,$-2$,$1$,$3$,$4$}
\tkzTabVar{+/$2$,-/$-1$,+/$3$,-/$1$,+/$2$}
\end{tikzpicture}
\end{questions}
\end{exo}

\partie{Extremums d'une fonction}

\begin{exo}[Lire les extremums]
\begin{questions}[label=\arabic*.]
  \item Maximum $\res{5}$, atteint en $1$ ; minimum $\res{-3}$, atteint en $4$
        ($f(-2)=1$ et $f(6)=0$ ne les dépassent pas).
  \item $\res{f'(x)\ge 0}$ sur $\intff{-2}{1}$, $\res{f'(x)\le 0}$ sur $\intff{1}{4}$,
        $\res{f'(x)\ge 0}$ sur $\intff{4}{6}$.
  \item En $1$ et en $4$, $f$ a un extremum local : la tangente y est horizontale, donc
        $f'(1)=\res{0}$ et $f'(4)=\res{0}$.
\end{questions}
\end{exo}

\begin{exo}[Extremums d'un polynôme]
\begin{questions}[label=\arabic*.]
  \item
\begin{calculs}
f'(x) &= 3x^{2}-12x+9 &&\\
f'(x) &= 3(x^{2}-4x+3) &&\justif{racines $1$ et $3$ ($\Delta=4$)}\\
f(1) &= 1-6+9+1=5 &&\\
f(3) &= 27-54+27+1=1 &&\\
f(5) &= 125-150+45+1=21
\end{calculs}
        Et $f(0)=1$.\\[2pt]
\begin{tikzpicture}[scale=0.8]
\tkzTabInit[lgt=1.6,espcl=2]{$x$/0.8,$f'(x)$/0.8,$f$/1.6}{$0$,$1$,$3$,$5$}
\tkzTabLine{,+,z,-,z,+,}
\tkzTabVar{-/$1$,+/$5$,-/$1$,+/$21$}
\end{tikzpicture}
  \item Maximum $\res{21}$, atteint en $5$ ; minimum $\res{1}$, atteint en $0$ et en $3$.
  \item Oui : $f'$ s'annule en $1$ en passant de $+$ à $-$ ; $f$ a un \res{maximum local
        égal à $5$} en $1$ (ce n'est pas le maximum sur $\intff{0}{5}$).
\end{questions}
\end{exo}

\begin{exo}[Vrai ou faux]
\begin{questions}[label=\arabic*.]
  \item \res{Faux.} Contre-exemple : $f(x)=x^{3}-6x^{2}+12x$ ;
        $f'(x)=3x^{2}-12x+12=3(x-2)^{2}$ s'annule en $2$ mais reste positive : $f$ est
        croissante, sans extremum en $2$. Il faut que $f'$ \emph{change de signe}.
  \item \res{Vrai.} $g'(x)=3x^{2}+1>0$ : $g$ est strictement croissante sur $\R$, sans
        extremum.
  \item \res{Vrai.}
\begin{calculs}
h'(x) &= 1-\frac{1}{x^{2}} &&\\
h'(x) &= \frac{x^{2}-1}{x^{2}} &&\\
h'(x) &= \frac{(x-1)(x+1)}{x^{2}} &&
\end{calculs}
        Sur $\intoo{0}{+\infty}$, $h'(x)$ a le signe de $x-1$ : $h$ décroît puis croît, et
        son minimum est $h(1)=1+1=2$.
\end{questions}
\end{exo}

\partie{Problèmes d'optimisation}

\begin{exo}[Un potager clôturé]
\begin{questions}[label=\arabic*.]
  \item Longueur $\times$ largeur $=100$, donc la longueur vaut $\dfrac{100}{x}$ et
\begin{calculs}
P(x) &= 2x+2\times\frac{100}{x} &&\\
P(x) &= 2x+\frac{200}{x}
\end{calculs}
  \item
\begin{calculs}
P'(x) &= 2-\frac{200}{x^{2}} &&\\
P'(x) &= \frac{2x^{2}-200}{x^{2}} &&\justif{même dénominateur}\\
P'(x) &= \frac{2(x-10)(x+10)}{x^{2}} &&\justif{$x^{2}-100=(x-10)(x+10)$}
\end{calculs}
  \item Sur $\intoo{0}{+\infty}$, $x^{2}>0$ et $x+10>0$ : $P'(x)$ a le signe de $x-10$.
        $P(10)=20+20=40$.\\[2pt]
\begin{tikzpicture}[scale=0.8]
\tkzTabInit[lgt=1.6,espcl=2.2]{$x$/0.8,$P'(x)$/0.8,$P$/1.6}{$0$,$10$,$+\infty$}
\tkzTabLine{d,-,z,+,}
\tkzTabVar{D+/{},-/$40$,+/{}}
\end{tikzpicture}
  \item Un carré de $\res{10~\text{m}}$ de côté ; il faut $\res{40~\text{m}}$ de
        clôture.
\end{questions}
\end{exo}

\begin{exo}[Bénéfice d'un atelier de vélos]
\begin{questions}[label=\arabic*.]
  \item
\begin{calculs}
B'(x) &= -6x^{2}+180x-750 &&\\
-6(x-5)(x-25) &= -6(x^{2}-30x+125) &&\\
-6(x-5)(x-25) &= -6x^{2}+180x-750 &&\justif{on retrouve $B'(x)$}
\end{calculs}
  \item $B'(x)$ est du signe de $-6$, sauf entre les racines $5$ et $25$.
\begin{calculs}
B(0) &= -1000 &&\\
B(5) &= -250+2250-3750-1000=-2750 &&\\
B(25) &= -31\,250+56\,250-18\,750-1000=5250 &&\\
B(40) &= -128\,000+144\,000-30\,000-1000=-15\,000
\end{calculs}
\begin{tikzpicture}[scale=0.8]
\tkzTabInit[lgt=1.6,espcl=2.2]{$x$/0.8,$B'(x)$/0.8,$B$/1.6}{$0$,$5$,$25$,$40$}
\tkzTabLine{,-,z,+,z,-,}
\tkzTabVar{+/$-1000$,-/$-2750$,+/$5250$,-/$-15\,000$}
\end{tikzpicture}
  \item Le maximum sur $\intff{0}{40}$ est $B(25)=5250$ (plus grand que $B(0)=-1000$) :
        il faut vendre $\res{25}$ vélos, pour un bénéfice de $\res{5250~\text{\euro}}$.
\end{questions}
\end{exo}

\begin{exo}[Une boîte avec le moins de carton]
\begin{questions}[label=\arabic*.]
  \item Volume $=x^{2}\times y=4$, donc $\res{y=\dfrac{4}{x^{2}}}$.
  \item Le fond mesure $x^{2}$ ; les $4$ faces latérales mesurent $x\times y$ chacune.
\begin{calculs}
S(x) &= x^{2}+4\times x\times\frac{4}{x^{2}} &&\\
S(x) &= x^{2}+\frac{16}{x}
\end{calculs}
  \item
\begin{calculs}
S'(x) &= 2x-\frac{16}{x^{2}} &&\\
S'(x) &= \frac{2x^{3}-16}{x^{2}} &&\justif{même dénominateur}\\
S'(x) &= \frac{2(x^{3}-8)}{x^{2}}
\end{calculs}
        La fonction cube est croissante et $2^{3}=8$ : $x^{3}-8<0$ si $x<2$ et
        $x^{3}-8>0$ si $x>2$. Comme $x^{2}>0$, $S'(x)$ a le même signe.
  \item $S(2)=4+8=12$.\\[2pt]
\begin{tikzpicture}[scale=0.8]
\tkzTabInit[lgt=1.6,espcl=2.2]{$x$/0.8,$S'(x)$/0.8,$S$/1.6}{$0$,$2$,$+\infty$}
\tkzTabLine{d,-,z,+,}
\tkzTabVar{D+/{},-/$12$,+/{}}
\end{tikzpicture}

        Il faut le moins de carton ($\res{12~\text{dm}^{2}}$) pour une base de
        $\res{2~\text{dm}}$ de côté et une hauteur $y=\dfrac{4}{4}=\res{1~\text{dm}}$.
\end{questions}
\end{exo}

\partie{Étude complète d'une fonction}

\begin{exo}[Étude d'un polynôme]
\begin{questions}[label=\arabic*.]
  \item $f$ est un polynôme : elle est \res{définie et dérivable sur $\R$}, donc sur
        $\intff{-1}{3}$.
  \item
\begin{calculs}
f'(x) &= -3x^{2}+6x &&\\
f'(x) &= \res{-3x(x-2)} &&\justif{racines $0$ et $2$}
\end{calculs}
  \item $f'(x)$ est du signe de $a=-3$ (négatif) à l'extérieur des racines, positif entre
        elles.
\begin{calculs}
f(-1) &= 1+3-1=3 &&\\
f(0) &= -1 &&\\
f(2) &= -8+12-1=3 &&\\
f(3) &= -27+27-1=-1
\end{calculs}
\begin{tikzpicture}[scale=0.8]
\tkzTabInit[lgt=1.6,espcl=2]{$x$/0.8,$f'(x)$/0.8,$f$/1.6}{$-1$,$0$,$2$,$3$}
\tkzTabLine{,-,z,+,z,-,}
\tkzTabVar{+/$3$,-/$-1$,+/$3$,-/$-1$}
\end{tikzpicture}
  \item Minimum local $\res{-1}$ en $0$ et maximum local $\res{3}$ en $2$ ($f'$ s'annule
        en changeant de signe). Sur $\intff{-1}{3}$ : maximum $\res{3}$ (atteint en $-1$
        et en $2$), minimum $\res{-1}$ (atteint en $0$ et en $3$).
  \item $f(1)=-1+3-1=1$ et $f'(1)=-3+6=3$.
\begin{calculs}
y &= 3(x-1)+1 &&\\
y &= \res{3x-2}
\end{calculs}
  \item
\begin{center}
\renewcommand{\arraystretch}{1.3}
\begin{tabular}{|c|c|c|c|c|c|c|c|c|c|}
\hline
$x$ & $-1$ & $-0{,}5$ & $0$ & $0{,}5$ & $1$ & $1{,}5$ & $2$ & $2{,}5$ & $3$\\ \hline
$f(x)$ & $3$ & $-0{,}125$ & $-1$ & $-0{,}375$ & $1$ & $2{,}375$ & $3$ & $2{,}125$ & $-1$\\ \hline
\end{tabular}

\smallskip
\begin{tikzpicture}[x=0.8cm,y=0.8cm]
  \repere{-2}{-2}{4}{4}
  \gradx{-1,1,2,3} \grady{-1,1,2,3}
  \draw[coulCourbe,line width=1.1pt,domain=-1:3,samples=100]
    plot (\x,{-\x*\x*\x+3*\x*\x-1});
  \draw[coulExo,line width=0.9pt] (0.2,-1.4) -- (1.8,3.4);
  \path[coulCourbe] (0,-1) node[croix]{} (2,3) node[croix]{} (1,1) node[croix]{};
  \node[coulExo,font=\footnotesize] at (1.75,3.75) {$T$};
  \node[coulCourbe,font=\footnotesize] at (3.4,-0.6) {$\Cf$};
\end{tikzpicture}
\end{center}
\end{questions}
\rem{En $1$, la courbe traverse sa tangente : ce n'est pas une erreur de tracé.}
\end{exo}

\begin{exo}[Étude d'un quotient défini sur $\R$]
\begin{questions}[label=\arabic*.]
  \item $x^{2}+2\ge 2>0$ : le dénominateur ne s'annule jamais, $f$ est \res{définie sur
        $\R$}.
  \item $u(x)=2x+1$, $u'(x)=2$ ; $v(x)=x^{2}+2$, $v'(x)=2x$.
\begin{calculs}
f'(x) &= \frac{2(x^{2}+2)-(2x+1)\times 2x}{(x^{2}+2)^{2}} &&\\
f'(x) &= \frac{2x^{2}+4-4x^{2}-2x}{(x^{2}+2)^{2}} &&\\
f'(x) &= \frac{-2x^{2}-2x+4}{(x^{2}+2)^{2}} &&\\
f'(x) &= \frac{-2(x+2)(x-1)}{(x^{2}+2)^{2}} &&\justif{$x^{2}+x-2$ a pour racines $-2$ et $1$}
\end{calculs}
  \item Le dénominateur est positif : $f'(x)$ a le signe de $-2x^{2}-2x+4$, négatif à
        l'extérieur des racines ($a=-2$), positif entre elles.
\begin{calculs}
f(-2) &= \frac{-3}{6}=-\frac{1}{2} &&\\
f(1) &= \frac{3}{3}=1
\end{calculs}
\begin{tikzpicture}[scale=0.8]
\tkzTabInit[lgt=1.6,espcl=2.2]{$x$/0.8,$f'(x)$/0.8,$f$/1.6}{$-\infty$,$-2$,$1$,$+\infty$}
\tkzTabLine{,-,z,+,z,-,}
\tkzTabVar{+/{},-/$-\frac{1}{2}$,+/$1$,-/{}}
\end{tikzpicture}
  \item Minimum local $\res{-\frac{1}{2}}$ en $-2$ ; maximum local $\res{1}$ en $1$.
  \item $f(0)=\frac{1}{2}$ et $f'(0)=\frac{4}{4}=1$, donc $\res{y=x+\frac{1}{2}}$.
  \item Valeurs arrondies au centième :
\begin{center}
\renewcommand{\arraystretch}{1.3}
\begin{tabular}{|c|c|c|c|c|c|c|c|c|c|}
\hline
$x$ & $-4$ & $-3$ & $-2$ & $-1$ & $0$ & $1$ & $2$ & $3$ & $4$\\ \hline
$f(x)$ & $-0{,}39$ & $-0{,}45$ & $-0{,}5$ & $-0{,}33$ & $0{,}5$ & $1$ & $0{,}83$ & $0{,}64$ & $0{,}5$\\ \hline
\end{tabular}

\smallskip
\begin{tikzpicture}[x=1cm,y=1cm]
  \repere{-5}{-2}{5}{2}
  \gradx{-4,-3,-2,-1,1,2,3,4} \grady{-1,1}
  \draw[coulCourbe,line width=1.1pt,domain=-4:4,samples=120]
    plot (\x,{(2*\x+1)/(\x*\x+2)});
  \draw[coulExo,line width=0.9pt] (-1.5,-1) -- (1.4,1.9);
  \path[coulCourbe] (-2,-0.5) node[croix]{} (1,1) node[croix]{} (0,0.5) node[croix]{};
  \node[coulExo,font=\footnotesize] at (1.6,1.75) {$T$};
  \node[coulCourbe,font=\footnotesize] at (3.8,0.85) {$\Cf$};
\end{tikzpicture}
\end{center}
\end{questions}
\end{exo}

\begin{exo}[Étude d'un quotient avec une valeur interdite]
\begin{questions}[label=\arabic*.]
  \item Le dénominateur s'annule en $2$ : $f$ est définie sur
        $\res{\R\setminus\{2\}}$.
  \item $u(x)=x^{2}-3$, $u'(x)=2x$ ; $v(x)=x-2$, $v'(x)=1$.
\begin{calculs}
f'(x) &= \frac{2x(x-2)-(x^{2}-3)\times 1}{(x-2)^{2}} &&\\
f'(x) &= \frac{2x^{2}-4x-x^{2}+3}{(x-2)^{2}} &&\\
f'(x) &= \frac{x^{2}-4x+3}{(x-2)^{2}} &&\\
f'(x) &= \frac{(x-1)(x-3)}{(x-2)^{2}} &&\justif{racines $1$ et $3$}
\end{calculs}
  \item $(x-2)^{2}>0$ : $f'(x)$ a le signe de $x^{2}-4x+3$, positif à l'extérieur des
        racines, négatif entre elles. La double barre marque la valeur interdite $2$.
\begin{calculs}
f(1) &= \frac{-2}{-1}=2 &&\\
f(3) &= \frac{6}{1}=6
\end{calculs}
\begin{tikzpicture}[scale=0.8]
\tkzTabInit[lgt=1.6,espcl=2]{$x$/0.8,$f'(x)$/0.8,$f$/1.6}{$-\infty$,$1$,$2$,$3$,$+\infty$}
\tkzTabLine{,+,z,-,d,-,z,+,}
\tkzTabVar{-/{},+/$2$,-D+/{}/{},-/$6$,+/{}}
\end{tikzpicture}
  \item Maximum local $\res{2}$ en $1$ ; minimum local $\res{6}$ en $3$. Le maximum
        local est plus petit que le minimum local : c'est possible, car la valeur
        interdite $2$ sépare les deux morceaux de la courbe.
  \item $f(0)=\frac{-3}{-2}=\frac{3}{2}$ et $f'(0)=\frac{(-1)(-3)}{4}=\frac{3}{4}$.
\begin{calculs}
y &= \frac{3}{4}(x-0)+\frac{3}{2} &&\\
y &= \res{\frac{3}{4}x+\frac{3}{2}}
\end{calculs}
  \item Valeurs arrondies au centième :
\begin{center}
\renewcommand{\arraystretch}{1.3}
\begin{tabular}{|c|c|c|c|c|c|c|}
\hline
$x$ & $-3$ & $-2$ & $-1$ & $0$ & $1$ & $1{,}5$\\ \hline
$f(x)$ & $-1{,}2$ & $-0{,}25$ & $0{,}67$ & $1{,}5$ & $2$ & $1{,}5$\\ \hline
\end{tabular}
\quad
\begin{tabular}{|c|c|c|c|c|c|c|}
\hline
$x$ & $2{,}5$ & $3$ & $4$ & $5$ & $6$ & $7$\\ \hline
$f(x)$ & $6{,}5$ & $6$ & $6{,}5$ & $7{,}33$ & $8{,}25$ & $9{,}2$\\ \hline
\end{tabular}

\smallskip
\begin{tikzpicture}[x=0.5cm,y=0.5cm]
  \repere{-4}{-2}{8}{10}
  \gradx{-3,-2,-1,1,2,3,4,5,6,7} \grady{-1,1,2,3,4,5,6,7,8,9}
  \draw[coulCourbe,line width=1.1pt,domain=-3:1.5,samples=80]
    plot (\x,{(\x*\x-3)/(\x-2)});
  \draw[coulCourbe,line width=1.1pt,domain=2.5:7,samples=80]
    plot (\x,{(\x*\x-3)/(\x-2)});
  \draw[coulExo,line width=0.9pt] (-3,-0.75) -- (2,3);
  \path[coulCourbe] (1,2) node[croix]{} (3,6) node[croix]{} (0,1.5) node[croix]{};
  \node[coulExo,font=\footnotesize] at (2.4,3.4) {$T$};
  \node[coulCourbe,font=\footnotesize] at (6.4,9.6) {$\Cf$};
\end{tikzpicture}
\end{center}
\end{questions}
\rem{On ne relie jamais les deux morceaux de la courbe : $f$ n'existe pas en $2$.}
\end{exo}

\partie{Problèmes concrets}

\begin{exo}[Le vol d'une montgolfière]
\begin{questions}[label=\arabic*.]
  \item $f(0)=\dfrac{6}{1}=6$ : le village est à $\res{600~\text{m}}$ d'altitude.
  \item $u(t)=8t+6$, $u'(t)=8$ ; $v(t)=t^{2}+1$, $v'(t)=2t$.
\begin{calculs}
f'(t) &= \frac{8(t^{2}+1)-(8t+6)\times 2t}{(t^{2}+1)^{2}} &&\\
f'(t) &= \frac{8t^{2}+8-16t^{2}-12t}{(t^{2}+1)^{2}} &&\\
f'(t) &= \frac{-8t^{2}-12t+8}{(t^{2}+1)^{2}} &&\\
f'(t) &= \frac{-4(2t^{2}+3t-2)}{(t^{2}+1)^{2}} &&\justif{on factorise par $-4$}\\
f'(t) &= \frac{-4(2t-1)(t+2)}{(t^{2}+1)^{2}} &&\justif{$(2t-1)(t+2)=2t^{2}+3t-2$}
\end{calculs}
  \item Sur $\intff{0}{5}$, $(t^{2}+1)^{2}>0$ et $t+2>0$ : $f'(t)$ a le signe de
        $-4(2t-1)$, positif pour $t<0{,}5$, négatif après.
\begin{calculs}
f(0{,}5) &= \frac{4+6}{1{,}25}=8 &&\\
f(5) &= \frac{46}{26}=\frac{23}{13} &&\justif{environ $1{,}8$}
\end{calculs}
\begin{tikzpicture}[scale=0.8]
\tkzTabInit[lgt=1.6,espcl=2.2]{$t$/0.8,$f'(t)$/0.8,$f$/1.6}{$0$,$0{,}5$,$5$}
\tkzTabLine{,+,z,-,}
\tkzTabVar{-/$6$,+/$8$,-/$\frac{23}{13}$}
\end{tikzpicture}
  \item $0{,}5$~h $=30$~min : la montgolfière commence à descendre à $\res{11~\text{h}~20}$,
        à $\res{800~\text{m}}$ d'altitude.
  \item On résout $f(t)=6$ avec $t>0$ :
\begin{calculs}
8t+6 &= 6(t^{2}+1) &&\justif{on multiplie par $t^{2}+1>0$}\\
6t^{2}-8t &= 0 &&\\
2t(3t-4) &= 0 &&
\end{calculs}
        $t=\dfrac{4}{3}$~h, soit $1$~h~$20$~min : elle revient à l'altitude du village à
        $\res{12~\text{h}~10}$.
\end{questions}
\end{exo}

\begin{exo}[Des bactéries et un antibiotique]
\begin{questions}[label=\arabic*.]
  \item $N(0)=10$ : il y a $\res{10\,000}$ bactéries au départ.
  \item
\begin{calculs}
N'(t) &= 3t^{2}-24t+36 &&\\
3(t-2)(t-6) &= 3(t^{2}-8t+12) &&\\
3(t-2)(t-6) &= 3t^{2}-24t+36 &&\justif{on retrouve $N'(t)$}
\end{calculs}
  \item $N'(t)$ est du signe de $3$ à l'extérieur des racines $2$ et $6$, négatif entre
        elles.
\begin{calculs}
N(2) &= 8-48+72+10=42 &&\\
N(6) &= 216-432+216+10=10 &&\\
N(8) &= 512-768+288+10=42
\end{calculs}
\begin{tikzpicture}[scale=0.8]
\tkzTabInit[lgt=1.6,espcl=2.2]{$t$/0.8,$N'(t)$/0.8,$N$/1.6}{$0$,$2$,$6$,$8$}
\tkzTabLine{,+,z,-,z,+,}
\tkzTabVar{-/$10$,+/$42$,-/$10$,+/$42$}
\end{tikzpicture}
  \item $N'(0)=36$ : au départ, la population augmente de $\res{36\,000}$ bactéries par
        heure. $N'(4)=3\times 2\times(-2)=-12$ : à $t=4$, elle diminue de
        $\res{12\,000}$ bactéries par heure.
  \item
\begin{calculs}
t(t-6)^{2} &= t(t^{2}-12t+36) &&\\
t(t-6)^{2} &= t^{3}-12t^{2}+36t &&\justif{c'est $N(t)-10$}
\end{calculs}
        Sur $\intff{0}{8}$, $t\ge 0$ et $(t-6)^{2}\ge 0$, donc $N(t)-10\ge 0$ :
        \res{non}, le nombre de bactéries ne descend jamais sous $10\,000$ (il y revient
        à $t=6$).
\end{questions}
\end{exo}

\begin{exo}[Le coût moyen d'un artisan]
\begin{questions}[label=\arabic*.]
  \item $M(10)=10+40=50$ : pour $10$ lampes, une lampe coûte en moyenne
        $\res{50~\text{\euro}}$.
  \item
\begin{calculs}
M'(x) &= 1-\frac{400}{x^{2}} &&\\
M'(x) &= \frac{x^{2}-400}{x^{2}} &&\justif{même dénominateur}\\
M'(x) &= \frac{(x-20)(x+20)}{x^{2}} &&\justif{$x^{2}-400=(x-20)(x+20)$}
\end{calculs}
  \item Sur $\intff{5}{50}$, $x^{2}>0$ et $x+20>0$ : $M'(x)$ a le signe de $x-20$.
\begin{calculs}
M(5) &= 5+80=85 &&\\
M(20) &= 20+20=40 &&\\
M(50) &= 50+8=58
\end{calculs}
\begin{tikzpicture}[scale=0.8]
\tkzTabInit[lgt=1.6,espcl=2.2]{$x$/0.8,$M'(x)$/0.8,$M$/1.6}{$5$,$20$,$50$}
\tkzTabLine{,-,z,+,}
\tkzTabVar{+/$85$,-/$40$,+/$58$}
\end{tikzpicture}

        Il faut fabriquer $\res{20}$ lampes ; le coût moyen est alors de
        $\res{40~\text{\euro}}$.
  \item On résout $M(x)\le 50$ :
\begin{calculs}
x+\frac{400}{x} &\le 50 &&\\
x^{2}+400 &\le 50x &&\justif{on multiplie par $x>0$}\\
x^{2}-50x+400 &\le 0
\end{calculs}
\begin{calculs}
\Delta &= 2500-1600=900 &&\\
x_{1} &= \frac{50-30}{2}=10 &&\\
x_{2} &= \frac{50+30}{2}=40
\end{calculs}
        Le trinôme est négatif entre ses racines : $M(x)\le 50$ pour $10\le x\le 40$.
        L'affirmation est \res{vraie}.
  \item $M'(10)=\dfrac{(-10)\times 30}{100}=-3$ : la tangente a pour coefficient
        directeur $-3$, et la droite $-2$. L'affirmation est \res{fausse}.
\end{questions}
\end{exo}

\end{document}
